% =========================================================================
% UNIT 2: EXPLORING TWO-VARIABLE DATA
% =========================================================================
\chapter{Unit 2: Exploring Two-Variable Data}

\section{Unit Overview \& Exam Weight}
Unit 2 accounts for \textbf{5\% to 7\%} of the AP Statistics Exam. Despite its modest percentage weight, regression analysis, correlation interpretation, residual evaluation, and computer output reading appear on virtually every single AP Exam, frequently as a high-scoring Free-Response Question (FRQ).

\subsection*{Key Unit Vocabulary}
\begin{description}[leftmargin=1.5em, style=nextline]
  \item[Explanatory Variable ($x$):] The independent variable hypothesized to explain or predict changes in the response variable.
  \item[Response Variable ($y$):] The dependent variable whose values are predicted or measured as an outcome.
  \item[Scatterplot:] A two-dimensional plot displaying the relationship between two quantitative variables measured on the same individuals.
  \item[Correlation Coefficient ($r$):] Measures the strength and direction of the linear relationship between two quantitative variables ($-1 \le r \le 1$).
  \item[Least-Squares Regression Line (LSRL):] The unique line that minimizes the sum of squared vertical residuals: $\min \sum (y - \hat{y})^2$.
  \item[Residual ($e$):] The directed difference between the observed value and the predicted value: $\text{Residual} = y - \hat{y}$.
  \item[Coefficient of Determination ($r^2$):] The percentage of total variability in $y$ explained by the linear relationship with $x$.
  \item[Extrapolation:] Predicting $y$ using $x$-values outside the domain of the collected data. Highly unreliable!
  \item[High Leverage Point:] An observation with an extreme $x$-value relative to other observations.
  \item[Influential Point:] An observation whose removal substantially changes the slope, $y$-intercept, or correlation.
\end{description}

\section{Core Concept Lessons}

\begin{lessonbox}[Lesson 2.1: Describing Scatterplots with D-O-F-S]
Whenever an exam question asks you to describe the relationship shown in a scatterplot, write complete sentences addressing all four components in context:
\begin{enumerate}[leftmargin=1.5em]
  \item \textbf{D --- Direction:} State whether the relationship is \textbf{Positive} (as $x$ increases, $y$ tends to increase) or \textbf{Negative} (as $x$ increases, $y$ tends to decrease).
  \item \textbf{O --- Outliers:} Identify any individual points that deviate substantially from the overall linear trend.
  \item \textbf{F --- Form:} State whether the pattern is \textbf{Linear} or \textbf{Non-linear} (curved).
  \item \textbf{S --- Strength:} Characterize the scatter as \textbf{Strong}, \textbf{Moderate}, or \textbf{Weak} based on how closely points cluster around a central trend line.
\end{enumerate}
\end{lessonbox}

\begin{lessonbox}[Lesson 2.2: The Correlation Coefficient ($r$)]
\begin{equation}
  r = \frac{1}{n-1}\sum_{i=1}^n \left(\frac{x_i - \bar{x}}{s_x}\right)\left(\frac{y_i - \bar{y}}{s_y}\right)
\end{equation}
\textbf{Essential Properties of $r$:}
\begin{itemize}[leftmargin=1.2em]
  \item $r$ is strictly bounded: $-1 \le r \le 1$.
  \item $r$ has \textbf{NO units of measurement}. Standardizing $x$ and $y$ makes $r$ invariant to changes in units (e.g., converting inches to centimeters does not change $r$).
  \item Switching $x$ and $y$ leaves $r$ completely unchanged!
  \item $r$ measures \textbf{LINEAR relationships only}! A scatterplot can have a perfect quadratic relationship ($y = x^2$) and still have $r = 0$.
  \item $r$ is \textbf{non-resistant}; an isolated outlier can dramatically alter its value.
  \item \textbf{Correlation does NOT imply causation!} An observed association may be driven by confounding variables.
\end{itemize}
\end{lessonbox}

\begin{lessonbox}[Lesson 2.3: The Least-Squares Regression Line (LSRL)]
\begin{equation}
  \hat{y} = b_0 + b_1 x \quad \text{or} \quad \hat{y} = a + bx
\end{equation}
\begin{equation}
  \text{Slope: } b_1 = r\left(\frac{s_y}{s_x}\right), \qquad y\text{-Intercept: } b_0 = \bar{y} - b_1 \bar{x}
\end{equation}
\textbf{Key Geometric Properties:}
\begin{itemize}[leftmargin=1.2em]
  \item The LSRL \textbf{always} passes through the point of means: $(\bar{x}, \bar{y})$.
  \item The sum of the vertical residuals is always mathematically equal to zero: $\sum (y - \hat{y}) = 0$.
  \item The sign of the slope $b_1$ is identical to the sign of the correlation coefficient $r$.
\end{itemize}
\end{lessonbox}

\begin{frqrubricbox}[Score-4 Interpretation Templates: Slope \& $y$-Intercept]
\textbf{Slope ($b_1$):}
\textit{"For each additional 1 [unit of $x$] increase in [explanatory variable], the predicted [response variable] increases/decreases by approximately $|b_1|$ [units of $y$]."}
\textit{(Note: You MUST include the word "predicted" or "estimated" to earn full credit!)}

\textbf{$y$-Intercept ($b_0$):}
\textit{"When the [explanatory variable] is 0 [units of $x$], the predicted [response variable] is $b_0$ [units of $y$]."}
\textit{(Note: Comment if $x=0$ makes physical sense in the real-world context).}
\end{frqrubricbox}

\begin{lessonbox}[Lesson 2.4: Residuals \& Residual Plot Analysis]
\begin{equation}
  \text{Residual } e = y - \hat{y} = \text{Actual } y - \text{Predicted } y
\end{equation}
\textbf{The Linear Appropriateness Test:}
A linear regression model is appropriate if and only if the residual plot exhibits a \textbf{random scatter of points around the zero line with roughly constant variance and no curved pattern}.
\begin{itemize}[leftmargin=1.2em]
  \item If the residual plot exhibits a curved (U-shaped or inverted U) pattern, the relationship between $x$ and $y$ is non-linear, and a linear model is \textbf{inappropriate}!
  \item A positive residual ($y > \hat{y}$) means the actual value was \textit{above} prediction (underestimate).
  \item A negative residual ($y < \hat{y}$) means the actual value was \textit{below} prediction (overestimate).
\end{itemize}
\end{lessonbox}

\begin{lessonbox}[Lesson 2.5: Coefficient of Determination ($r^2$) \& Residual SD ($s$)]
\begin{frqrubricbox}[Score-4 Template: $r^2$ and $s$ Interpretation]
\textbf{Coefficient of Determination ($r^2$):}
\textit{"Approximately $[r^2 \times 100\%]$ of the variability in [response variable $y$] is accounted for by the linear relationship with [explanatory variable $x$]."}

\textbf{Standard Deviation of Residuals ($s$):}
\begin{equation}
  s = \sqrt{\frac{\sum (y - \hat{y})^2}{n - 2}}
\end{equation}
\textit{"When using the least-squares regression line to predict [response variable $y$] from [explanatory variable $x$], our predictions will typically be off by approximately $s$ [units of $y$]."}
\end{frqrubricbox}
\end{lessonbox}

\begin{lessonbox}[Lesson 2.6: High Leverage, Outliers, \& Influential Points]
\begin{itemize}[leftmargin=1.2em]
  \item \textbf{High Leverage Point:} Has an $x$-value substantially higher or lower than the mean $\bar{x}$. Acts as an anchor point; if it aligns with the overall trend, it inflates $r$ without altering slope.
  \item \textbf{Outlier in Regression:} An observation with a very large residual (far above or below the regression line).
  \item \textbf{Influential Point:} An observation whose removal markedly changes the regression slope, $y$-intercept, or correlation. High leverage points that deviate from the linear trend are typically highly influential on slope!
\end{itemize}
\end{lessonbox}

\begin{lessonbox}[Lesson 2.7: Reading Regression Computer Output]
Computer printouts (from Minitab, R, or JMP) always follow this standardized grid:
\begin{center}
\begin{adjustbox}{max width=\linewidth}
\begin{tabular}{lcccc}
\toprule
\textbf{Predictor} & \textbf{Coef} & \textbf{SE Coef} & \textbf{T} & \textbf{P} \\
\midrule
\textbf{Constant} & $b_0$ & $\text{SE}(b_0)$ & $t$ & $p$ \\
\textbf{Variable ($x$)} & $b_1$ & $\text{SE}(b_1)$ & $t$ & $p$ \\
\bottomrule
\end{tabular}
\end{adjustbox}
\end{center}
$S = s$ (Standard deviation of residuals), $R\text{-sq} = r^2$.
To obtain $r$ from $R\text{-sq}$, compute $\sqrt{r^2}$ and \textbf{match the sign of the slope ($b_1$)}!
\end{lessonbox}

\section{Worked Step-by-Step Examples}

\subsection*{Example 2.1: Interpreting Computer Output \& Calculating Residuals}
\textbf{Problem:} A biologist investigates the relationship between ambient water temperature ($x$, in $^\circ\text{C}$) and the larval development time of a species of fish ($y$, in days). Regression output for 12 tanks is displayed below:
\begin{center}
\begin{adjustbox}{max width=\linewidth}
\begin{tabular}{lcccc}
\toprule
\textbf{Predictor} & \textbf{Coef} & \textbf{SE Coef} & \textbf{T} & \textbf{P} \\
\midrule
\textbf{Constant} & 48.650 & 2.140 & 22.73 & 0.000 \\
\textbf{Temp} & -1.420 & 0.095 & -14.95 & 0.000 \\
\bottomrule
\end{tabular}
\end{adjustbox}
\end{center}
$S = 1.15\text{ days}, \quad R\text{-sq} = 95.7\%, \quad R\text{-sq(adj)} = 95.3\%$.\\
(a) Write the equation of the least-squares regression line.\\
(b) Interpret the slope of the regression line in context.\\
(c) Calculate and interpret the correlation coefficient $r$.\\
(d) One tank at $18^\circ\text{C}$ had an observed larval development time of $21.5\text{ days}$. Calculate the residual.

\textbf{Solution:}
\begin{itemize}[leftmargin=1.2em]
  \item \textbf{Part (a):} $\widehat{\text{Development Time}} = 48.650 - 1.420(\text{Temperature})$.
  \item \textbf{Part (b):} \textit{"For each additional $1^\circ\text{C}$ increase in ambient water temperature, the predicted larval development time decreases by approximately $1.420\text{ days}$."}
  \item \textbf{Part (c):} $r = -\sqrt{0.957} \approx -0.978$. (We assign the negative sign because the slope $b_1 = -1.420$ is negative).
  \textit{"There is a very strong, negative linear relationship between water temperature and larval development time."}
  \item \textbf{Part (d):} Predicted development time for $x = 18$:
  \[ \hat{y} = 48.650 - 1.420(18) = 48.650 - 25.56 = 23.09\text{ days} \]
  $\text{Residual } e = y - \hat{y} = 21.5 - 23.09 = -1.59\text{ days}$.\\
  \textit{"The actual development time was $1.59\text{ days}$ shorter than predicted by the linear regression model."}
\end{itemize}

\section{TI-84 Plus CE Calculator Playbook}

\begin{calculatorbox}[TI-84 Playbook: Linear Regression \& Residual Plots]
\begin{enumerate}[leftmargin=1.2em]
  \item \textbf{Diagnostic ON:} Press \texttt{[2nd] [0]} (\texttt{CATALOG}) $\rightarrow$ Scroll to \texttt{DiagnosticOn} $\rightarrow$ Press \texttt{[ENTER]} twice. (Crucial to display $r$ and $r^2$!).
  \item Enter $x$-data into \texttt{L1} and $y$-data into \texttt{L2}.
  \item Press \texttt{[STAT]} $\rightarrow$ \texttt{CALC} $\rightarrow$ Select \texttt{8:LinReg(a+bx)}.
  \item Set \texttt{Xlist: L1}, \texttt{Ylist: L2}, \texttt{Store RegEQ: Y1} (press \texttt{[VARS] $\rightarrow$ Y-VARS $\rightarrow$ 1:Function $\rightarrow$ 1:Y1}). Press \texttt{Calculate}.
  \item \textbf{To plot residuals:} In \texttt{STAT PLOT}, set \texttt{Ylist: RESID} (press \texttt{[2nd] [STAT]} $\rightarrow$ select \texttt{RESID}). Press \texttt{[ZOOM] [9:ZoomStat]}.
\end{enumerate}
\end{calculatorbox}

\section{Comprehensive Practice Problem Set}

\subsection*{Part I: 20 Multiple-Choice Questions}

\begin{enumerate}[leftmargin=1.5em]
  \item If the correlation between two variables $x$ and $y$ is $r = 0.85$, what is the correlation if $x$ is converted from feet to meters ($1\text{ ft} = 0.3048\text{ m}$)?
  \begin{enumerate}[label=(\Alph*)]
    \item $0.85 \times 0.3048$ \item $0.85 / 0.3048$ \item 0.85 \item $-0.85$ \item 0.7225
  \end{enumerate}
  \textbf{Answer: (C).} Changing measurement units does not alter the correlation coefficient $r$.

  \item A regression line has equation $\hat{y} = 12 - 3.5x$. What is the predicted change in $y$ for each 2-unit increase in $x$?
  \begin{enumerate}[label=(\Alph*)]
    \item Decreases by 3.5 \item Decreases by 7.0 \item Increases by 7.0 \item Increases by 12 \item Decreases by 12
  \end{enumerate}
  \textbf{Answer: (B).} $2 \times (-3.5) = -7.0$.

  \item Which of the following indicates that a linear model is appropriate?
  \begin{enumerate}[label=(\Alph*)]
    \item The correlation $r$ is positive \item The residual plot shows a clear curved pattern \item The residual plot displays random scatter around the horizontal zero line \item $r^2$ is exactly 0.50 \item The slope is greater than 1
  \end{enumerate}
  \textbf{Answer: (C).} Random scatter without systematic curvature confirms appropriateness.

  \item An observation has an actual value of $y = 45$ and a predicted value of $\hat{y} = 52$. The residual is:
  \begin{enumerate}[label=(\Alph*)]
    \item $-7$ \item $+7$ \item 45 \item 52 \item 0.865
  \end{enumerate}
  \textbf{Answer: (A).} $\text{Residual} = y - \hat{y} = 45 - 52 = -7$.

  \item If $r^2 = 0.64$ and the slope of the regression line is negative, what is $r$?
  \begin{enumerate}[label=(\Alph*)]
    \item 0.64 \item 0.80 \item $-0.80$ \item $-0.64$ \item 0.4096
  \end{enumerate}
  \textbf{Answer: (C).} $r = -\sqrt{0.64} = -0.80$.

  \item The point of means $(\bar{x}, \bar{y})$:
  \begin{enumerate}[label=(\Alph*)]
    \item Always lies on the least-squares regression line \item Lies on the line only if $r = 1$ \item Is always an outlier \item Has a non-zero residual \item Determines the direction of curvature
  \end{enumerate}
  \textbf{Answer: (A).} Mathematically, $\bar{y} = b_0 + b_1\bar{x}$.

  \item What does $r^2 = 0.81$ mean in the context of predicting fuel efficiency ($y$) from car weight ($x$)?
  \begin{enumerate}[label=(\Alph*)]
    \item 81\% of cars have high fuel efficiency.
    \item 81\% of the variability in fuel efficiency is accounted for by the linear model with car weight.
    \item The correlation is 0.81.
    \item The slope of the line is 0.81.
    \item Predictions will be correct 81\% of the time.
  \end{enumerate}
  \textbf{Answer: (B).} Standard definition of the coefficient of determination.

  \item An influential point in regression is best described as an observation that:
  \begin{enumerate}[label=(\Alph*)]
    \item Has a residual of zero.
    \item Substantially alters the slope, intercept, or correlation when removed.
    \item Is always an error in data collection.
    \item Falls exactly on the mean of $x$.
    \item Has a positive $z$-score.
  \end{enumerate}
  \textbf{Answer: (B).} Defining property of an influential observation.

  \item Using a model $\hat{y} = 5 + 2x$ established for $x \in [10, 50]$ to predict $y$ when $x = 120$ is an example of:
  \begin{enumerate}[label=(\Alph*)]
    \item Interpolation \item Extrapolation \item Randomization \item Confounding \item Residual analysis
  \end{enumerate}
  \textbf{Answer: (B).} Predicting beyond the range of explanatory data is extrapolation.

  \item If $s_x = 4, s_y = 6,$ and $r = 0.75$, what is the slope $b_1$ of the LSRL?
  \begin{enumerate}[label=(\Alph*)]
    \item 0.75 \item 1.125 \item 0.50 \item 1.50 \item 0.888
  \end{enumerate}
  \textbf{Answer: (B).} $b_1 = r(s_y / s_x) = 0.75(6/4) = 0.75(1.5) = 1.125$.
\end{enumerate}

\begin{enumerate}[resume, leftmargin=1.5em]
  \item Which of the following correlations indicates the strongest linear relationship?
  \begin{enumerate}[label=(\Alph*)]
    \item $r = 0.78$ \item $r = 0.00$ \item $r = -0.89$ \item $r = 0.85$ \item $r = -0.12$
  \end{enumerate}
  \textbf{Answer: (C).} Strength is determined by $|r|$; $|-0.89| = 0.89$ is the largest magnitude.

  \item If the standard deviation of residuals is $s = 3.2\text{ cm}$, this means:
  \begin{enumerate}[label=(\Alph*)]
    \item All residuals are less than 3.2 cm.
    \item Predictions of $y$ typically deviate from actual values by about 3.2 cm.
    \item The slope is 3.2.
    \item $r = 0.32$.
    \item The mean residual is 3.2 cm.
  \end{enumerate}
  \textbf{Answer: (B).} Standard interpretation of $s$.

  \item If the residuals for a model sum to a non-zero value, what is the most likely explanation?
  \begin{enumerate}[label=(\Alph*)]
    \item The relationship is non-linear.
    \item Rounding error in calculations.
    \item The line used is not the least-squares regression line.
    \item $r$ was negative.
    \item There are outliers.
  \end{enumerate}
  \textbf{Answer: (C).} The true LSRL strictly satisfies $\sum e = 0$; any other line will not.

  \item High leverage points in a scatterplot always have:
  \begin{enumerate}[label=(\Alph*)]
    \item Large residuals \item Unusual $x$-values \item Negative slopes \item $r > 0.90$ \item Low variance
  \end{enumerate}
  \textbf{Answer: (B).} Leverage is governed purely by the distance of $x$ from $\bar{x}$.

  \item When two variables have a strong non-linear relationship of the form $y = a b^x$, what transformation linearizes the data?
  \begin{enumerate}[label=(\Alph*)]
    \item Plotting $y$ vs. $x^2$ \item Plotting $\log(y)$ vs. $x$ \item Plotting $\log(y)$ vs. $\log(x)$ \item Plotting $1/y$ vs. $1/x$ \item Plotting $\sqrt{y}$ vs. $\sqrt{x}$
  \end{enumerate}
  \textbf{Answer: (B).} An exponential model is linearized by taking the logarithm of $y$ only ($\log y = \log a + (\log b)x$).

  \item A power model $y = a x^b$ is linearized by plotting:
  \begin{enumerate}[label=(\Alph*)]
    \item $\log(y)$ vs. $x$ \item $\log(y)$ vs. $\log(x)$ \item $y$ vs. $\log(x)$ \item $y^2$ vs. $x$ \item $\sqrt{y}$ vs. $x$
  \end{enumerate}
  \textbf{Answer: (B).} Power models linearize when both variables are transformed logarithmically.

  \item In a regression analysis, the explanatory variable is speed (mph) and response is stopping distance (ft). The slope is 3.5. What are the units of the slope?
  \begin{enumerate}[label=(\Alph*)]
    \item mph \item ft \item ft per mph \item mph per ft \item Dimensionless
  \end{enumerate}
  \textbf{Answer: (C).} Slope units are always $[y\text{-units}] / [x\text{-units}] = \text{ft per mph}$.

  \item If $r = 0$, what does this conclude about the relationship between $x$ and $y$?
  \begin{enumerate}[label=(\Alph*)]
    \item There is no relationship of any kind.
    \item There is no linear relationship between $x$ and $y$.
    \item All residuals are zero.
    \item The slope is 1.
    \item The data is non-random.
  \end{enumerate}
  \textbf{Answer: (B).} A strong curved relationship can still have $r = 0$; $r$ detects only linear trends.

  \item A researcher observes a strong positive correlation ($r = 0.92$) between ice cream sales and drowning deaths. The most plausible explanation is:
  \begin{enumerate}[label=(\Alph*)]
    \item Eating ice cream causes drowning.
    \item Drowning causes people to buy ice cream.
    \item Both variables are influenced by a common confounding variable: outdoor summer temperature.
    \item The correlation was miscalculated.
    \item Extrapolation error.
  \end{enumerate}
  \textbf{Answer: (C).} Classic lurking/confounding variable scenario.

  \item A regression line has $\hat{y} = 10 + 2x$. If an observation at $x = 5$ has an actual value of $y = 23$, the point lies:
  \begin{enumerate}[label=(\Alph*)]
    \item Exactly on the line \item 3 units below the line \item 3 units above the line \item 23 units above the line \item 10 units below the line
  \end{enumerate}
  \textbf{Answer: (C).} $\hat{y} = 10 + 2(5) = 20$. Actual $y = 23 \implies \text{Residual} = 23 - 20 = +3$.
\end{enumerate}

\subsection*{Part II: Free-Response Practice Problem}

\begin{workbookbox}[FRQ 2.1: Regression Modeling \& Residual Evaluation]
\textbf{Problem:} An environmental engineer investigates the relationship between wind speed ($x$, in mph) and electricity generation ($y$, in kilowatts) of a wind turbine. The regression line is $\widehat{\text{Power}} = -15.4 + 4.8(\text{Wind})$. $r^2 = 0.88, s = 6.2\text{ kW}$.
(a) Interpret the slope $4.8$ in context.\\
(b) Interpret $r^2 = 0.88$ in context.\\
(c) At a wind speed of $20\text{ mph}$, the turbine produced $85\text{ kW}$. Calculate the residual and state whether the model overpredicted or underpredicted power.

\textbf{Grader-Approved Score-4 Solution:}
\begin{itemize}[leftmargin=1.2em]
  \item \textbf{Part (a):} \textit{"For each additional 1 mph increase in wind speed, the predicted electricity generation increases by approximately 4.8 kilowatts."}
  \item \textbf{Part (b):} \textit{"Approximately 88\% of the variability in electricity generation is accounted for by the linear relationship with wind speed."}
  \item \textbf{Part (c):} Predicted: $\hat{y} = -15.4 + 4.8(20) = -15.4 + 96.0 = 80.6\text{ kW}$.\\
  $\text{Residual} = y - \hat{y} = 85 - 80.6 = \mathbf{+4.4\text{ kW}}$.\\
  Because the residual is positive, the linear regression model \textbf{underpredicted} the actual electricity generation by 4.4 kW.
\end{itemize}
\end{workbookbox}


\begin{frqrubricbox}[FRQ 2.2: Regression Output, Residuals, and Outliers vs. High Leverage]
\textbf{Problem:} A real estate analyst investigates the relationship between living area ($x$, in hundreds of square feet) and selling price ($y$, in thousands of dollars) for 24 recently sold suburban homes. A computer regression package outputs:
\begin{center}
\begin{adjustbox}{max width=\linewidth}
\begin{tabular}{lcccc}
\toprule
\textbf{Predictor} & \textbf{Coef} & \textbf{SE Coef} & \textbf{T} & \textbf{P} \\
\midrule
Constant & 45.20 & 16.40 & 2.76 & 0.011 \\
Area & 12.80 & 1.60 & 8.00 & 0.000 \\
\bottomrule
\end{tabular}
\end{adjustbox}
\end{center}
$S = 18.50\text{ (\$1,000s)}, \quad R\text{-sq} = 74.4\%$.
\begin{enumerate}[label=(\alph*)]
  \item State the equation of the least-squares regression line and predict the selling price of a home with 2,500 square feet ($x = 25.0$).
  \item A home with 2,500 square feet sold for \$390,000 ($y = 390.0$). Calculate and interpret the residual.
  \item Distinguish between an outlier in regression and a point with high leverage. If a 6,000-sq-ft mansion is added to the dataset, would it be classified as high leverage? Explain.
  \item Explain why the analyst should not use this model to predict the price of a 7,500-square-foot commercial estate.
\end{enumerate}

\textbf{Grader-Approved Score-4 Solution \& Rubric:}
\begin{itemize}[leftmargin=1.2em]
  \item \textbf{Part (a):} $\widehat{\text{Price}} = 45.20 + 12.80(\text{Area})$.\\
  For $x = 25.0$: $\hat{y} = 45.20 + 12.80(25.0) = 45.20 + 320.00 = 365.20\text{ thousand dollars}$ (\$365,200).
  \item \textbf{Part (b):} $\text{Residual} = y - \hat{y} = 390.0 - 365.20 = +24.80\text{ thousand dollars}$ (\$24,800).\\
  \textit{Interpretation:} The actual selling price of the home was \$24,800 higher than the price predicted by the regression line based on its living area.
  \item \textbf{Part (c):} An outlier has an unusually large residual (far above or below the regression line). A point with high leverage has an extreme $x$-value (far from $\bar{x}$) with the potential to strongly steer the slope. A 6,000-sq-ft home ($x = 60.0$) is far beyond the typical sample range, giving it high leverage.
  \item \textbf{Part (d):} Predicting at 7,500 square feet represents dangerous \textbf{extrapolation}. We have no evidence that the linear relationship continues beyond the observed data range.
\end{itemize}
\end{frqrubricbox}

\begin{trapbox}[Unit 2 Top 5 AP Exam Pitfalls to Avoid]
\begin{enumerate}[leftmargin=1.2em]
  \item \textbf{Implying Causation from Correlation:} $r$ measures linear association, never cause and effect!
  \item \textbf{Claiming $r^2$ Means the Model is Correct:} $r^2$ can be high even if the underlying data is curved. Always inspect the residual plot!
  \item \textbf{Confusing Residual Direction:} Residual is Observed minus Predicted ($y - \hat{y}$), never $\hat{y} - y$.
  \item \textbf{Omitting Units in Slope Interpretation:} Always state "predicted increase in [units of $y$] per 1 [unit of $x$]".
  \item \textbf{Extrapolation Danger:} Never make predictions outside the domain of observed $x$-values.
\end{enumerate}
\end{trapbox}
\section{Unit 2 Master Summary}
\begin{lessonbox}[Unit 2 Essential Review Checklist]
\begin{itemize}[leftmargin=1.2em]
  \item \textbf{D-O-F-S:} Direction, Outliers, Form, Strength for scatterplots.
  \item \textbf{LSRL Formula:} $\hat{y} = b_0 + b_1 x$, where $b_1 = r(s_y/s_x)$ and $b_0 = \bar{y} - b_1\bar{x}$.
  \item \textbf{Residual:} $e = y - \hat{y}$ (Actual minus Predicted).
  \item \textbf{Residual Plot Appropriateness:} Random scatter around zero confirms a linear model is appropriate; any curved pattern indicates a non-linear relationship.
  \item \textbf{Interpretation Keywords:} Always use "predicted" or "estimated" for slope interpretations.
\end{itemize}
\end{lessonbox}
\clearpage
